%============================================================
%  Bd. 31 - Rechtecksbänder
%============================================================

\documentclass[main.tex]{subfiles}

\ifSubfilesClassLoaded{
  \BandSetupStandalone{B31}
}{
}

\title{Bd. 31 - Rechtecksbänder}
\author{Martin Kunze}
\date{}
\setcounter{file}{31}

\hypersetup{
  pdftitle={Bd. 31 - Rechtecksbänder},
  pdfauthor={Martin Kunze}
}

\begin{document}

\title{Bd. 31 - Rechtecksbänder}
\author{Martin Kunze}
\date{}
\hypersetup{pageanchor=false}
\maketitle
\hypersetup{pageanchor=true}
\setcounter{tocdepth}{3}
\tableofcontents
\setcounter{file}{31}
\renewcommand{\FormulaBandID}{31}

\chapter{Einführung}

Dieser Band führt Rechtecksbänder als besondere Halbgruppen ein. Die
Rechtecksidentität erzwingt bereits die Idempotenz. Anschließend werden das
Dreitermgesetz und die Produktdarstellung bewiesen. Die McLean-Zerlegung
liefert darauf aufbauend die strukturelle Grundlage für die Rekonstruktion
endlicher idempotenter Halbgruppen aus ihren Potenzhalbgruppen.

\chapter{Rechtecksbänder}

\section{Definition und grundlegende Eigenschaften}
\FormulaDefDeltaK[Begriff des Rechtecksbandes]{%
  \RectangularBandAlg{R}{\star}%
}{RectangularBandDef}{%
  \DeltaRow{Mengen}{R\dsep x\dsep y}
  \DeltaRow{\textbf{Axiome}}{}
  \DeltaPrem{Halbgruppe}{\SemigroupAlg{R}{\star}}
  \DeltaRow{Rechtecksidentität}{%
    x,y\in R\vdash (x\star y)\star x=x}
  \DeltaRow{\textbf{Neues Symbol}}{}
  \DeltaPrem{Rechtecksband}{%
    \RectangularBandAlg{R}{\star}}
}

\FormulaAxiomDeltaK[Zugriff auf die Halbgruppe eines Rechtecksbandes]{%
  \RectangularBandAlg{R}{\star}
  \vdash\SemigroupAlg{R}{\star}%
}{RectangularBandBaseAxiom}{%
  \DeltaRow{Mengen}{R}
}

\FormulaAxiomDeltaK[Zugriff auf die Rechtecksidentität]{%
  \RectangularBandAlg{R}{\star}\dsep x,y\in R
  \vdash (x\star y)\star x=x%
}{RectangularBandIdentityAxiom}{%
  \DeltaRow{Mengen}{R\dsep x\dsep y}
}

\FormulaThmDeltaK[Rechtecksbänder sind idempotente Halbgruppen]{%
  \RectangularBandAlg{R}{\star}
  \vdash\IdempotentSemigroupAlg{R}{\star}%
}{RectangularBandIsIdempotentSemigroup}{%
  \DeltaRow{Mengen}{R\dsep x}
  \DeltaRow{Funktionen}{\star}
}
\begin{tabproofwide}
  \proofstepwidestar[1]{\RectangularBandAlg{R}{\star}}{\rA}
  \proofstepwidestar[1]{\SemigroupAlg{R}{\star}}{%
    \FormulaRefAuto{RectangularBandBaseAxiom}[ax]{1}}
  \proofstepwidestar[3]{x\in R}{\rA}
  \proofstepwidestar[1,3]{x\star x\in R}{%
    \FormulaRefAuto{SemigroupClosureAxiom}[ax]{2,3}}
  \proofstepwidestar[1,3]{(x\star x)\star x=x}{%
    \FormulaRefAuto{RectangularBandIdentityAxiom}[ax]{1,3}}
  \proofstepwidestar[1,3]{%
    (x\star x)\star x=x\star(x\star x)}{%
    \FormulaRefAuto{SemigroupAssociativityAxiom}[ax]{2,3}}
  \proofstepwide[1,3]{x\star x}{=}{%
    \bigl((x\star x)\star x\bigr)\star x}{%
    \hspace*{1em}%
    \FormulaRefAuto{BinaryFunctionLeftArgumentCongruence}[thm]{%
      \FormulaRefAuto{a=b\vdash b=a}[thm]{5}}}
  \proofstepwide[1,3]{}{=}{%
    \bigl(x\star(x\star x)\bigr)\star x}{%
    \hspace*{1em}%
    \FormulaRefAuto{BinaryFunctionLeftArgumentCongruence}[thm]{6}}
  \proofstepwide[1,3]{}{=}{x}{%
    \FormulaRefAuto{RectangularBandIdentityAxiom}[ax]{1,3,4}}
  \proofstepwidestar[1,3]{x\star x=x}{%
    \rChain{7,9}}
  \proofstepwidestar[1]{\IdempotentSemigroupAlg{R}{\star}}{%
    \FormulaRefAuto{IdempotentSemigroupDef}[def]{2,10}}
\end{tabproofwide}

\FormulaThmDeltaK[Dreitermgesetz der Rechtecksbänder]{%
  \begin{aligned}[t]
    &\RectangularBandAlg{R}{\star}\dsep x,y,z\in R\vdash{}\\[-2pt]
    &(x\star y)\star z=x\star z
  \end{aligned}%
}{RectangularBandThreeTermLaw}{%
  \DeltaRow{Mengen}{R\dsep x\dsep y\dsep z}
  \DeltaRow{Funktionen}{\star}
}
\begin{tabproofwide}
  \proofstepwidestar[1]{\RectangularBandAlg{R}{\star}}{\rA}
  \proofstepwidestar[1]{\SemigroupAlg{R}{\star}}{%
    \FormulaRefAuto{RectangularBandBaseAxiom}[ax]{1}}
  \proofstepwidestar[3]{x,y,z\in R}{\rA}
  \proofstepwidestar[1,3]{(z\star x)\star z=z}{%
    \FormulaRefAuto{RectangularBandIdentityAxiom}[ax]{1,3}}
  \proofstepwidestar[1,3]{y\star z\in R}{%
    \FormulaRefAuto{SemigroupClosureAxiom}[ax]{2,3}}
  \proofstepwidestar[1,3]{%
    \bigl(x\star(y\star z)\bigr)\star x=x}{%
    \FormulaRefAuto{RectangularBandIdentityAxiom}[ax]{1,3,5}}
  \proofstepwide[1,3]{(x\star y)\star z}{=}{%
    (x\star y)\star\bigl((z\star x)\star z\bigr)}{%
    \hspace*{1em}%
    \FormulaRefAuto{a=b\vdash b=a}[thm]{%
      \FormulaRefAuto{BinaryFunctionRightArgumentCongruence}[thm]{4}}}
  \proofstepwide[1,3]{}{=}{%
    \bigl((x\star(y\star z))\star x\bigr)\star z}{%
    \hspace*{1em}%
    \FormulaRefAuto{SemigroupFiveFactorReassociation}[thm]{2,3}}
  \proofstepwide[1,3]{}{=}{x\star z}{%
    \FormulaRefAuto{BinaryFunctionLeftArgumentCongruence}[thm]{6}}
  \proofstepwidestar[1,3]{(x\star y)\star z=x\star z}{%
    \rChain{7,9}}
\end{tabproofwide}

\FormulaThmDeltaK[Dreitermcharakterisierung der Rechtecksbänder]{%
  \begin{aligned}[t]
    &\RectangularBandAlg{R}{\star}\eqvdash\\[-2pt]
    &\left(
      \IdempotentSemigroupAlg{R}{\star}\land
      \forall x,y,z\in R\,
        \bigl((x\star y)\star z=x\star z\bigr)
      \right)
  \end{aligned}%
}{RectangularBandThreeTermCharacterization}{%
  \DeltaRow{Mengen}{R\dsep u\dsep v\dsep w\dsep x\dsep y\dsep z}
  \DeltaRow{Funktionen}{\star}
}

\ProofWideReasonsNearFormula
\begin{tabproofsplitwide}
  \proofpartwide{\(\vdash\)}
    \proofstepwidestar[1]{\RectangularBandAlg{R}{\star}}{\rA}
    \proofstepwidestar[1]{\IdempotentSemigroupAlg{R}{\star}}{%
      \FormulaRefAuto{RectangularBandIsIdempotentSemigroup}[thm]{1}}
    \proofstepwidestar[3]{x,y,z\in R}{\rA}
    \proofstepwidestar[1,3]{(x\star y)\star z=x\star z}{%
      \FormulaRefAuto{RectangularBandThreeTermLaw}[thm]{1,3}}
    \proofstepwidestarbreak[1]{%
      \forall x,y,z\in R\,
        \bigl((x\star y)\star z=x\star z\bigr)}{%
      \rUI{\rRI{3,4}}}
    \proofstepwidestarbreak[1]{%
      \IdempotentSemigroupAlg{R}{\star}\land
      \forall x,y,z\in R\,
        \bigl((x\star y)\star z=x\star z\bigr)}{%
      \rAI{2,5}}
  \closeproofpartwide

  \proofpartwide{\(\dashv\)}
    \proofstepwidestarbreakkeep[1]{%
      \IdempotentSemigroupAlg{R}{\star}\land
      \forall u,v,w\in R\,
        \bigl((u\star v)\star w=u\star w\bigr)}{\rA}
    \proofstepwidestar[1]{\IdempotentSemigroupAlg{R}{\star}}{\rAEa{1}}
    \proofstepwidestarbreak[1]{%
      \forall u,v,w\in R\,
        \bigl((u\star v)\star w=u\star w\bigr)}{\rAEb{1}}
    \proofstepwidestar[1]{\SemigroupAlg{R}{\star}}{%
      \FormulaRefAuto{IdempotentSemigroupBaseAxiom}[ax]{2}}
    \proofstepwidestar[5]{x,y\in R}{\rA}
    \proofstepwidestar[1,5]{x\star x=x}{%
      \FormulaRefAuto{IdempotentSemigroupIdempotenceAxiom}[ax]{2,5}}
    \proofstepwidestarinline[1,5]{%
      (x\star y)\star x=x\star x}{%
      \FormulaRefAuto{\forall x \in A\,(P(x)),\, y \in A
        \vdash P(y)}[thm]{%
        \FormulaRefAuto{\forall x \in A\,(P(x)),\, y \in A
          \vdash P(y)}[thm]{%
          \FormulaRefAuto{\forall x \in A\,(P(x)),\, y \in A
            \vdash P(y)}[thm]{3,5},5},5}}
    \proofstepwide[1,5]{(x\star y)\star x}{=}{x}{%
      \rChain{7,6}}
    \proofstepwidestar[1]{\RectangularBandAlg{R}{\star}}{%
      \FormulaRefAuto{RectangularBandDef}[def]{4,8}}
  \closeproofpartwide
\end{tabproofsplitwide}
\ProofWideReasonsAtRight

\FormulaThmDeltaK[Koordinatensatz f\"ur Rechtecksb\"ander]{%
  \begin{aligned}[t]
    &R\neq\varnothing\vdash{}\\[-2pt]
    &\RectangularBandAlg{R}{\star}
      \leftrightarrow
      \forall e\in R\,
        \SemigroupCoordinateIso{R}{\star}{e}
  \end{aligned}%
}{RectangularBandProductRepresentation}{%
  \DeltaRow{Mengen}{%
    R\dsep e\dsep p\dsep q\dsep r\dsep s\dsep
    x\dsep y\dsep z}
  \DeltaRow{Funktionen}{\star\dsep
    \SemigroupCoordinateMap{R}{\star}{e}\dsep
    \SemigroupCoordinateProductOp{R}{\star}{e}}
}

Für \(e\in R\) bezeichnet
\(\SemigroupCoordinateIso{R}{\star}{e}\) den in Band~28 eingeführten
\(e\)-Koordinatenisomorphismus.  Seine konkrete Koordinatenabbildung ist
\[
  \SemigroupCoordinateMap{R}{\star}{e}(x)
  =(x\star e,e\star x)\qquad(x\in R),
\]
sein Ziel ist
\(\SemigroupCoordinateCarrier{R}{\star}{e}\), und dort gilt die
Koordinatenregel
\[
  (p,q)\SemigroupCoordinateProductOp{R}{\star}{e}(r,s)=(p,s)
  \qquad
  \bigl((p,q),(r,s)\in
    \SemigroupCoordinateCarrier{R}{\star}{e}\bigr).       \tag{25.1}
\]
Damit ist die rechte Seite eine axiomatisch benannte Produktdarstellung
und nicht länger eine erneute Aufzählung ihrer Einzelbedingungen.

\ProofWideReasonsNearFormula
\begin{tabproofsplitwide}
  \proofpartwideindR[Rechtecksband \(\Rightarrow\)
    Koordinatenisomorphismen]{%
    \begin{aligned}[t]
      &R\neq\varnothing\dsep
        \RectangularBandAlg{R}{\star}\vdash{}\\[-2pt]
      &\forall e\in R\,
        \SemigroupCoordinateIso{R}{\star}{e}
    \end{aligned}
  }{%
    \begin{aligned}[t]
      &R\neq\varnothing\dsep
        \RectangularBandAlg{R}{\star}\vdash{}\\[-2pt]
      &\forall e\in R\,
        \SemigroupCoordinateIso{R}{\star}{e}
    \end{aligned}
  }[RectangularBandRepresentationForward]
    \proofstepwidestar[1]{R\neq\varnothing}{\rA}
    \proofstepwidestar[2]{\RectangularBandAlg{R}{\star}}{\rA}
    \proofstepwidestar[2]{\SemigroupAlg{R}{\star}}{%
      \FormulaRefAuto{RectangularBandBaseAxiom}[ax]{2}}
    \proofstepwidestar[2]{\IdempotentSemigroupAlg{R}{\star}}{%
      \FormulaRefAuto{RectangularBandIsIdempotentSemigroup}[thm]{2}}
    \proofstepwidestar[5]{x,y,z\in R}{\rA}
    \proofstepwidestar[2,5]{(x\star y)\star z=x\star z}{%
      \FormulaRefAuto{RectangularBandThreeTermLaw}[thm]{2,5}}
    \proofstepwidestarbreak[2]{%
      \forall x,y,z\in R\,
        \bigl((x\star y)\star z=x\star z\bigr)}{%
      \rUI{\rRI{5,6}}}
    \proofstepwidestar[8]{e\in R}{\rA}
    \proofstepwidestar[9]{x\in R}{\rA}
    \proofstepwidestar[2,9]{x\star x=x}{%
      \FormulaRefAuto{IdempotentSemigroupIdempotenceAxiom}[ax]{4,9}}
    \proofstepwidestarbreak[2]{%
      \forall x\in R\,(x\star x=x)}{%
      \rUI{\rRI{9,10}}}
    \proofstepwidestar[2,8]{%
      \SemigroupCoordinateIso{R}{\star}{e}}{%
      \FormulaRefAuto{SemigroupCoordinateIsoCriterion}[thm]{3,8,11,7}}
    \proofstepwidestar[2]{%
      \forall e\in R\,
        \SemigroupCoordinateIso{R}{\star}{e}}{%
      \hspace*{1em}%
      \rUI{\rRI{8,12}}}
  \closeproofpartwideindR

  \proofpartwideindR[Koordinatenisomorphismen \(\Rightarrow\)
    Rechtecksband]{%
    \begin{aligned}[t]
      &R\neq\varnothing\dsep
        \forall e\in R\,
          \SemigroupCoordinateIso{R}{\star}{e}
        \vdash\RectangularBandAlg{R}{\star}
    \end{aligned}
  }{%
    \begin{aligned}[t]
      &R\neq\varnothing\dsep
        \forall e\in R\,
          \SemigroupCoordinateIso{R}{\star}{e}
        \vdash\RectangularBandAlg{R}{\star}
    \end{aligned}
  }[RectangularBandRepresentationBackward]
    \proofstepwidestar[1]{R\neq\varnothing}{\rA}
    \proofstepwidestar[2]{%
      \forall e\in R\,
        \SemigroupCoordinateIso{R}{\star}{e}}{%
      \hspace*{1em}%
      \rA}
    \proofstepwidestar[1]{\exists e\,(e\in R)}{%
      \FormulaRefAuto{A\neq\varnothing\eqvdash
        \exists x\,(x\in A)}[thm]{1}}
    \proofstepwidestar[4]{e\in R}{\rA}
    \proofstepwidestar[2,4]{%
      \SemigroupCoordinateIso{R}{\star}{e}}{%
      \FormulaRefAuto{\forall x \in A\,(P(x)),\, y \in A
        \vdash P(y)}[thm]{2,4}}
    \proofstepwidestar[2,4]{\SemigroupAlg{R}{\star}}{%
      \FormulaRefAuto{SemigroupCoordinateIsoSourceAxiom}[ax]{5}}
    \proofstepwidestarbreak[2,4]{%
      \forall x,y\in R\,
        \bigl((x\star y)\star x=x\bigr)}{%
      \FormulaRefAuto{SemigroupCoordinateIsoRectangularIdentity}[thm]{5}}
    \proofstepwidestar[2,4]{\RectangularBandAlg{R}{\star}}{%
      \FormulaRefAuto{RectangularBandDef}[def]{6,7}}
    \proofstepwidestar[1,2]{\RectangularBandAlg{R}{\star}}{%
      \rEE{3,4,8}}
  \closeproofpartwideindR

  \proofpartwide{Schluss}
    \proofstepwidestar[1]{R\neq\varnothing}{\rA}
    \proofstepwidestar[2]{\RectangularBandAlg{R}{\star}}{\rA}
    \proofstepwidestar[1,2]{%
      \forall e\in R\,
        \SemigroupCoordinateIso{R}{\star}{e}}{%
      \hspace*{.75em}%
      \FormulaRefAuto{RectangularBandRepresentationForward}[thm]{1,2}}
    \proofstepwidestarbreak[1]{%
      \RectangularBandAlg{R}{\star}\rightarrow
      \forall e\in R\,
        \SemigroupCoordinateIso{R}{\star}{e}}{\rRI{2,3}}
    \proofstepwidestar[5]{%
      \forall e\in R\,
        \SemigroupCoordinateIso{R}{\star}{e}}{\rA}
    \proofstepwidestar[1,5]{\RectangularBandAlg{R}{\star}}{%
      \FormulaRefAuto{RectangularBandRepresentationBackward}[thm]{1,5}}
    \proofstepwidestarbreak[1]{%
      \left(
        \forall e\in R\,
          \SemigroupCoordinateIso{R}{\star}{e}
      \right)\rightarrow
      \RectangularBandAlg{R}{\star}}{\rRI{5,6}}
    \proofstepwidestarbreakkeep[1]{%
      \RectangularBandAlg{R}{\star}
      \leftrightarrow
      \forall e\in R\,
        \SemigroupCoordinateIso{R}{\star}{e}}{\rLRI{4,7}}
  \closeproofpartwide
  \enlargethispage{\baselineskip}
\end{tabproofsplitwide}
\ProofWideReasonsAtRight


\chapter{Rekonstruktion aus der Potenzhalbgruppe}

In diesem Kapitel schreiben wir die Halbgruppenoperation als Juxtaposition
und kürzen \(\PowNE{S}\) durch \(\mathcal P_+(S)\) ab.  Für
\(a\in S\) steht \(\Psi(a)\) stets für \(\Psi(\{a\})\); diese Menge
muss nicht einelementig sein.  Gerade diese Unterscheidung ist wesentlich:
Bei Linksnullhalbgruppen können Isomorphismen der Potenzhalbgruppen
Einermengen auf größere Mengen abbilden.

Ziel ist die für die endliche Potenzhalbgruppenvermutung benötigte
Richtung
\[
  \mathcal P_+(S)\cong\mathcal P_+(T)
  \quad\Longrightarrow\quad S\cong T
\]
für endliche idempotente Halbgruppen.  Der Beweis folgt der
Rekonstruktionsidee von
\href{https://doi.org/10.1080/00927872.2017.1319474}%
  {Baomin Yu, Xianzhong Zhao und Aiping Gan,
  \emph{Global determinism of idempotent semigroups}, 2018},
trennt aber den endlichen Kardinalitätsschritt ausdrücklich von den
strukturellen Argumenten.
Bis zum Hauptsatz seien die betrachteten Grundhalbgruppen nichtleer; der
Hauptsatz behandelt den Leerfall vor Anwendung der Bandzerlegung gesondert.

\section{Bandzerlegung}

Fixiere f\"ur ein nichtleeres Rechtecksband ein \(e\in R\) und setze
\[
 I=\{x\star e\mid x\in R\},\qquad
 J=\{e\star x\mid x\in R\}.
\]
Dann ist \(\SemigroupCoordinateCarrier{R}{\star}{e}=I\times J\).
Entlang des Isomorphismus
\(\SemigroupCoordinateMap{R}{\star}{e}\) identifizieren wir \(R\) mit
diesem Koordinatentr\"ager; sein Koordinatenprodukt aus (25.1) wird wie
vereinbart durch Juxtaposition geschrieben.  In dieser Darstellung ist der Fall
\(|J|=1\) ein
Linksnullband, der Fall \(|I|=1\) ein Rechtsnullband; sind beide Faktoren
einelementig, so ist \(R\) trivial.  Sind \(|I|,|J|>1\), nennen wir
\(R\) ein \emph{echtes Rechtecksband}.

Wir verwenden den Struktursatz für Bänder: Jede idempotente Halbgruppe
besitzt eine maximale Halbverbandszerlegung
\[
  S=[Y;S_\alpha],\qquad
  S=\mathop{\dot\bigcup}_{\alpha\in Y}S_\alpha,           \tag{25.2}
\]
wobei \(Y\) eine kommutative idempotente Halbgruppe, jedes
\(S_\alpha\) ein Rechtecksband und
\[
  S_\alpha S_\beta\subseteq S_{\alpha\beta}              \tag{25.3}
\]
ist.  Die Ordnung auf \(Y\) ist durch
\(\beta\leq\alpha\iff\alpha\beta=\beta\) gegeben.
Konkret ist (25.2) der Quotient nach der kleinsten Kongruenz, deren
Quotient kommutativ ist; ihre Klassen sind genau die maximalen
Rechtecksbänder.  Idempotenz liefert die Idempotenz des Quotienten,
und (25.3) folgt aus der Kongruenzeigenschaft.  Dies ist die
McLean-Zerlegung eines Bandes.

Der folgende Komponentensatz ist der strukturelle Ausgangspunkt des
Rekonstruktionsbeweises.  Er ist wichtig genug, um seine beiden Aussagen
getrennt festzuhalten.

\begin{quote}
\textbf{Komponentensatz.}
Seien
\(S=[Y;S_\alpha]\) und \(T=[Z;T_\gamma]\) Bandzerlegungen und sei
\[
  \Psi\colon\mathcal P_+(S)\longrightarrow\mathcal P_+(T)
\]
ein Halbgruppenisomorphismus.  Dann existiert ein
Halbverbandsisomorphismus \(\theta\colon Y\to Z\) mit
\[
  \Psi\bigl[\mathcal P_+(S_\alpha)\bigr]
     =\mathcal P_+(T_{\theta(\alpha)})                    \tag{25.4}
\]
für jedes \(\alpha\in Y\).  Sind
\(a\in S_\alpha\), \(b\in S_\beta\) und
\(\alpha\nleq\beta\), so ist außerdem
\[
  \Psi(aba)                                                \tag{25.5}
\]
eine Einermenge.
\end{quote}

Die Komponentenerkennung (25.4) ist Satz~3.3 und die
Sandwicherkennung (25.5) ist Proposition~3.8(ii) in
\href{https://doi.org/10.1080/00927872.2015.1087006}%
  {Aiping Gan, Xianzhong Zhao und Yong Shao,
  \emph{Globals of Idempotent Semigroups},
  \emph{Communications in Algebra} 44 (2016), Nr.~9, 3743--3766}.
Dieser frühere Artikel liefert hier nur die beiden Vorresultate; seine
globale Bestimmtheitsaussage betrifft normale Bänder, nicht alle
idempotenten Halbgruppen.  Die Aussage für die ganze Klasse wird erst in der
oben verlinkten späteren Arbeit von Yu--Zhao--Gan als Satz~4.9 formuliert;
die hier benötigte endliche Fassung wird unten eigenständig bewiesen.
Die dort für \(\alpha>\beta\) formulierte Sandwichaussage liefert
(25.5) auch unter der schwächeren Voraussetzung
\(\alpha\nleq\beta\): Dann ist
\(c=aba\in S_{\alpha\beta}\), es gilt
\(\alpha>\alpha\beta\) und \(aca=c\), sodass die Aussage auf das Paar
\(a,c\) anwendbar ist.
Sie beruhen auf der maximalen Halbverbandszerlegung der Potenzhalbgruppe;
insbesondere setzen sie die zu rekonstruierenden Einermengen nicht voraus.
Damit entsteht kein Zirkelschluss mit dem folgenden Beweis.

\paragraph{Endliche echte Rechtecksbänder erkennen Einermengen.}
Seien \(R=I\times J\) und \(R'=I'\times J'\) endliche echte
Rechtecksbänder und
\(\Phi\colon\mathcal P_+(R)\to\mathcal P_+(R')\) ein
Isomorphismus.  Dann bildet \(\Phi\) jede Einermenge auf eine Einermenge
ab und induziert einen Isomorphismus \(R\to R'\).

\begingroup
\let\proofstepwidestar\proofstepwidestarforcebreak
\let\proofstepwidestarbreak\proofstepwidestarforcebreak
\ProofWideReasonsNearFormula
\begin{tabproofwide}
  \proofstepwidestar[]{%
    \varnothing\neq A,B\subseteq I\times J
    \vdash AB=\pi_I(A)\times\pi_J(B)\quad\text{\normalfont(25.6)}}{%
    \text{\normalfont Koordinatenregel (25.1)}}
  \proofstepwidestar[]{%
    \begin{aligned}[t]
      A\equiv_R B&\iff
        (\forall C\in\mathcal P_+(R))\ AC=BC,\\[-2pt]
      A\equiv_L B&\iff
        (\forall C\in\mathcal P_+(R))\ CA=CB
    \end{aligned}}{%
    \text{\normalfont rein multiplikative Definition}}
  \proofstepwidestar[]{%
    \begin{aligned}[t]
      A\equiv_RB&\iff\pi_I(A)=\pi_I(B),\\[-2pt]
      A\equiv_LB&\iff\pi_J(A)=\pi_J(B)
    \end{aligned}
    \quad\text{\normalfont(25.7)}}{%
    \text{\normalfont Schritte 1--2; Einsetzen der Produktformel (25.6)}}
  \proofstepwidestarbreak[]{%
    \begin{aligned}[t]
      &K\subseteq I\dsep |K|=k\vdash{}\\[-2pt]
      &\bigl|\{A\in\mathcal P_+(R)\mid\pi_I(A)=K\}\bigr|
        =(2^{|J|}-1)^k
    \end{aligned}}{%
    \text{\normalfont je Zeile eine nichtleere Teilmenge wählen}}
  \proofstepwidestarbreak[]{%
    \begin{aligned}[t]
      &A\text{ liegt in einer kleinsten }\equiv_R\text{-Klasse}
        \iff |\pi_I(A)|=1,\\[-2pt]
      &A\text{ liegt in einer kleinsten }\equiv_L\text{-Klasse}
        \iff |\pi_J(A)|=1
    \end{aligned}}{%
    \text{\normalfont Schritte 3--4; \(|J|,|I|>1\); zweiter Schluss dual}}
  \proofstepwidestar[]{%
    \Phi\text{ erhält }\equiv_R,\equiv_L
      \text{ und die endlichen Klassengrößen}}{%
    \text{\normalfont Isomorphie und Bijektivität von \(\Phi\)}}
  \proofstepwidestarbreak[]{%
    |\pi_I(A)|=|\pi_J(A)|=1
      \iff
    |\pi_{I'}(\Phi(A))|=|\pi_{J'}(\Phi(A))|=1}{%
    \text{\normalfont Schritte 5--6; multiplikative Erkennung}}
  \proofstepwidestar[]{%
    A\in\mathcal P_+(I\times J)\dsep
    |\pi_I(A)|=|\pi_J(A)|=1
    \iff \exists x\in I\times J\,(A=\{x\})}{%
    \text{\normalfont zwei einelementige Projektionen}}
  \proofstepwidestar[]{%
    \Phi\text{ bildet Einermengen bijektiv auf Einermengen ab}}{%
    \text{\normalfont Schritte 7--8; Charakterisierung und Bijektivität}}
  \proofstepwidestarbreak[]{%
    \Phi(\{x\}\{y\})=\Phi(\{xy\})
    \quad(x,y\in R)}{%
    \text{\normalfont Multiplikativität von \(\Phi\)}}
  \proofstepwidestar[]{%
    \Phi\text{ induziert einen Isomorphismus }R\longrightarrow R'}{%
    \text{\normalfont Schritte 9--10; Einschränkung auf die Einermengenschichten}}
\end{tabproofwide}
\endgroup

Ist eine Komponente ein Links- beziehungsweise Rechtsnullband, so ist auch
die durch (25.4) zugeordnete Komponente vom gleichen Typ: In ihrer
Potenzhalbgruppe gilt dann \(AB=A\) beziehungsweise \(AB=B\), und die
Gleichung für Einermengen überträgt sich auf die Grundkomponente.  Eine
triviale Komponente bleibt trivial, weil (25.4) eine Bijektion ihrer
nichtleeren Potenzmengen liefert.  Somit stimmen die vier Komponententypen
unter \(\theta\) überein.

\section{Transport innerhalb der Komponenten}

Von nun an seien \(S,T\) endlich, und \(\theta\) und \(\Psi\) seien wie
im Komponentensatz.  Für eine nichtleere Teilmenge \(A\subseteq S\) und
\(U\subseteq T\) setzen wir
\[
\begin{aligned}
 \widehat\Psi(A)
   &\coloneqq\bigcup_{\varnothing\neq C\subseteq A}\Psi(C),\\
 \widehat{\Psi^{-1}}(U)
   &\coloneqq\bigcup_{\varnothing\neq V\subseteq U}\Psi^{-1}(V).
\end{aligned}                                                   \tag{25.8}
\]
Beide Operatoren sind monoton.  Ferner gelten die Inklusionen
\[
 \Psi(A)\subseteq\widehat\Psi(A),
 \qquad
 A\subseteq\widehat{\Psi^{-1}}(\Psi(A)),
\]
sowie ihre dualen Entsprechungen.

\paragraph{Sandwichrechnung.}
Sind \(a\in S_\alpha\), \(b\in S_\beta\) und
\(\alpha\nleq\beta\), so gilt
\begin{align}
&\widehat{\Psi^{-1}}(\Psi(a))\,
 \widehat{\Psi^{-1}}(\Psi(b))\,
 \widehat{\Psi^{-1}}(\Psi(a))=\{aba\},                  \tag{25.9}\\
&\widehat\Psi\!\left(\widehat{\Psi^{-1}}(\Psi(a))\right)
 \widehat\Psi\!\left(\widehat{\Psi^{-1}}(\Psi(b))\right)
 \widehat\Psi\!\left(\widehat{\Psi^{-1}}(\Psi(a))\right)
 =\Psi(aba).                                               \tag{25.10}
\end{align}

\begingroup
\let\proofstepwidestar\proofstepwidestarforcebreak
\let\proofstepwidestarbreak\proofstepwidestarforcebreak
\ProofWideReasonsNearFormula
\begin{tabproofwide}
  \proofstepwidestarbreak[]{%
    \begin{aligned}[t]
      &x_1\in\widehat{\Psi^{-1}}(\Psi(a))\dsep
       y\in\widehat{\Psi^{-1}}(\Psi(b))\dsep
       x_2\in\widehat{\Psi^{-1}}(\Psi(a))\\[-2pt]
      &\vdash\exists U_1,U_2\in\mathcal P_+(\Psi(a))\,
       \exists V\in\mathcal P_+(\Psi(b))\;\bigl(
       x_i\in\Psi^{-1}(U_i)\ (i=1,2)\land
       y\in\Psi^{-1}(V)\bigr)
    \end{aligned}}{%
    \text{\normalfont Definition (25.8)}}
  \proofstepwidestarbreak[]{%
    \varnothing\neq U_1VU_2
      \subseteq\Psi(a)\Psi(b)\Psi(a)=\Psi(aba)}{%
    \text{\normalfont Teilmengenprodukt und Multiplikativität von \(\Psi\)}}
  \proofstepwidestar[]{\Psi(aba)\text{ ist eine Einermenge}}{%
    \text{\normalfont Sandwicherkennung (25.5)}}
  \proofstepwidestar[]{U_1VU_2=\Psi(aba)}{%
    \text{\normalfont Schritte 2--3; nichtleere Teilmenge einer Einermenge}}
  \proofstepwidestarbreak[]{%
    \Psi^{-1}(U_1)\Psi^{-1}(V)\Psi^{-1}(U_2)
      =\Psi^{-1}(U_1VU_2)=\{aba\}}{%
    \text{\normalfont Schritt 4; Multiplikativität von \(\Psi^{-1}\)}}
  \proofstepwidestar[]{%
    \widehat{\Psi^{-1}}(\Psi(a))\,
    \widehat{\Psi^{-1}}(\Psi(b))\,
    \widehat{\Psi^{-1}}(\Psi(a))\subseteq\{aba\}}{%
    \text{\normalfont Schritte 1 und 5; beliebige Elemente der Faktoren}}
  \proofstepwidestar[]{%
    a\in\widehat{\Psi^{-1}}(\Psi(a))\dsep
    b\in\widehat{\Psi^{-1}}(\Psi(b))}{%
    \text{\normalfont Inklusionen nach (25.8)}}
  \proofstepwidestar[]{%
    \widehat{\Psi^{-1}}(\Psi(a))\,
    \widehat{\Psi^{-1}}(\Psi(b))\,
    \widehat{\Psi^{-1}}(\Psi(a))=\{aba\}}{%
    \text{\normalfont Schritte 6--7; beide Inklusionen; dies ist (25.9)}}
  \proofstepwidestarbreak[]{%
    \widehat\Psi\!\left(\widehat{\Psi^{-1}}(\Psi(a))\right)
    \widehat\Psi\!\left(\widehat{\Psi^{-1}}(\Psi(b))\right)
    \widehat\Psi\!\left(\widehat{\Psi^{-1}}(\Psi(a))\right)
    =\Psi(aba)}{%
    \text{\normalfont Schritt 8; \(\Psi\) anwenden und nach (25.8) vereinigen}}
\end{tabproofwide}
\endgroup

Aus (25.9) folgt zunächst noch folgende Mengenform.  Sind
\(a\in S_\alpha\), \(\alpha>\beta\),
\(B\in\mathcal P_+(S_\beta)\) und
\(A_1,A_2\) nichtleere Teilmengen von
\(\widehat{\Psi^{-1}}(\Psi(a))\), so gilt
\[
  A_1BA_2=aBa.                                             \tag{25.10c}
\]
Denn für jedes \(b\in B\) ist (25.9), auf die entsprechenden
Teilmengen eingeschränkt, genau
\(A_1\{b\}A_2=\{aba\}\); die Vereinigung über \(b\in B\) ergibt
(25.10c).

Insbesondere folgen für \(\alpha>\beta\) die beiden Formen
\[
 \Psi(a)t\Psi(a)\text{ ist einelementig}
 \quad(a\in S_\alpha,\ t\in T_{\theta(\beta)}),           \tag{25.10a}
\]
und
\[
 s\Psi(b)s\text{ ist einelementig}
 \quad(s\in T_{\theta(\alpha)},\ b\in S_\beta).          \tag{25.10b}
\]
Für (25.10a) wähle \(r\in\Psi(a)\) und setze
\(A_r=\Psi^{-1}(\{r\})\),
\(B_t=\Psi^{-1}(\{t\})\).  Dann ist
\(A_r\subseteq\widehat{\Psi^{-1}}(\Psi(a))\), und (25.10c) liefert
\[
 A_rB_tA_r=aB_ta.
\]
Nach Anwenden von \(\Psi\) ist daher
\(\Psi(a)t\Psi(a)=\{rtr\}\).  Für (25.10b) wenden wir (25.10a)
zunächst auf \(\Psi^{-1}\) an.  Dadurch ist
\(\Psi^{-1}(\{s\})\,b\,\Psi^{-1}(\{s\})\) einelementig, etwa gleich
\(\{cbc\}\) mit \(c\in\Psi^{-1}(\{s\})\).  Nach Anwenden von \(\Psi\)
ist \(s\Psi(b)s=\Psi(cbc)\), und diese Menge ist nach (25.5)
einelementig.

Für \(\alpha>\beta\) definieren wir auf \(S_\alpha\)
\[
 x\mathrel{\sigma_{\alpha,\beta}}y
 \iff
 (\forall b\in S_\beta)\ (xb=yb\ \land\ bx=by),          \tag{25.11}
\]
und für \(\gamma>\alpha\) auf \(S_\alpha\)
\[
 x\mathrel{\kappa_{\alpha,\gamma}}y
 \iff
 (\forall c\in S_\gamma)\ cxc=cyc.                      \tag{25.12}
\]
Beide Relationen sind Kongruenzen.  Für \(\sigma\) folgt dies unmittelbar
daraus, dass Multiplikation mit einem Element aus \(S_\alpha\) die
Testelemente in \(S_\beta\) belässt.  Sind nämlich
\(x\mathrel{\sigma_{\alpha,\beta}}y\),
\(d\in S_\alpha\) und \(b\in S_\beta\), so gilt
\[
\begin{aligned}
 (dx)b&=d(xb)=d(yb)=(dy)b,&
 b(dx)&=(bd)x=(bd)y=b(dy),\\
 (xd)b&=x(db)=y(db)=(yd)b,&
 b(xd)&=(bx)d=(by)d=b(yd),
\end{aligned}
\]
weil \(bd,db\in S_\beta\).  Für \(\kappa\) zeigt etwa
\[
 c(xd)c=cx(xc)dc=(cxc)dc=(cyc)dc=cy(yc)dc=c(yd)c
\]
für \(d\in S_\alpha\), \(c\in S_\gamma\) die Verträglichkeit mit
Rechtsmultiplikation; links gilt das duale Argument.  Dabei wurde in der
Rechteckskomponente \(S_\alpha\) die Identität \(uvw=uw\) benutzt.

Aus der Idempotenz erhält man außerdem die nützliche Charakterisierung
\[
 x\mathrel{\sigma_{\alpha,\beta}}y
 \iff
 (\forall b\in S_\beta)\ xbx=yby.                         \tag{25.13}
\]
Die Hinrichtung folgt aus
\(xbx=(xb)(bx)\).  Umgekehrt liefern
\(xb=(xbx)b\) und \(bx=b(xbx)\) die beiden Gleichungen in
(25.11).

Setze
\[
 \sigma_\alpha=\bigcap_{\beta<\alpha}\sigma_{\alpha,\beta},
 \qquad
 \kappa_\alpha=\bigcap_{\gamma>\alpha}\kappa_{\alpha,\gamma},
 \qquad
 \rho_\alpha=\sigma_\alpha\cap\kappa_\alpha,             \tag{25.14}
\]
wobei ein leerer Durchschnitt die Allrelation auf \(S_\alpha\) bedeutet.
Dann ist \(\rho_\alpha\) eine Kongruenz.  Auf den Komponenten von \(T\)
werden \(\sigma',\kappa',\rho'\) gleichartig definiert.

\begin{quote}
\textbf{Transportlemma.}
Für \(a,b\in S_\alpha\) gilt
\[
 a\mathrel{\rho_\alpha}b
 \iff
 (\forall s,t\in\Psi(a)\cup\Psi(b))\;
 s\mathrel{\rho'_{\theta(\alpha)}}t.                    \tag{25.15}
\]
Ist \(C=a\rho_\alpha\) und \(t\in\Psi(a)\), so bildet \(\Psi\)
\(\mathcal P_+(C)\) isomorph auf
\(\mathcal P_+(t\rho'_{\theta(\alpha)})\) ab.  Die Vorschrift
\[
 \tau_\alpha(a\rho_\alpha)
   =t\rho'_{\theta(\alpha)}\qquad(t\in\Psi(a))            \tag{25.16}
\]
definiert daher einen Isomorphismus
\[
 \tau_\alpha\colon S_\alpha/\rho_\alpha
    \longrightarrow T_{\theta(\alpha)}/\rho'_{\theta(\alpha)}.
                                                                  \tag{25.17}
\]
Schließlich gelten für \(a\in S_\alpha\), \(b\in S_\beta\):
\begin{align}
 \alpha,\beta>\alpha\beta:\quad&
 \tau_\alpha(a\rho_\alpha)\,
 \tau_\beta(b\rho_\beta)=\Psi(ab),                       \tag{25.18}\\[2mm]
 \alpha>\beta:\quad&
 \tau_\alpha(a\rho_\alpha)\,
 \tau_\beta(b\rho_\beta)\,
 \tau_\alpha(a\rho_\alpha)=\Psi(aba).                    \tag{25.19}
\end{align}
Die Produkte auf den linken Seiten sind jeweils Einermengen.
\end{quote}

\begingroup
\let\proofstepwidestar\proofstepwidestarforcebreak
\let\proofstepwidestarbreak\proofstepwidestarforcebreak
\ProofWideReasonsNearFormula
\begin{tabproofsplitwide}
  \proofpartwide{Transport der \(\sigma\)-Relation}
    \proofstepwidestar[1]{x\mathrel{\sigma_{\alpha,\beta}}y}{\rA}
    \proofstepwidestar[1]{%
      \forall b\in S_\beta\;(xbx=yby)}{%
      \text{\normalfont Charakterisierung (25.13)}}
    \proofstepwidestarbreak[1]{%
      s\mathrel{\sigma'_{\theta(\alpha),\theta(\beta)}}t
      \quad(s,t\in\Psi(x))}{%
      \text{\normalfont Schritt 2; (25.13) und Sandwichrechnung}}
    \proofstepwidestarbreak[1]{%
      s\mathrel{\sigma'_{\theta(\alpha),\theta(\beta)}}t
      \quad(s,t\in\Psi(y))}{%
      \text{\normalfont Schritt 2; (25.13) und Sandwichrechnung}}
    \proofstepwidestarbreak[]{%
      B'\in\mathcal P_+(T_{\theta(\beta)})\dsep
      B=\Psi^{-1}(B')
      \vdash B\in\mathcal P_+(S_\beta)}{%
      \text{\normalfont Komponententransport (25.4)}}
    \proofstepwidestarbreak[1]{%
      \Psi(x)B'=\Psi(y)B'\dsep
      B'\Psi(x)=B'\Psi(y)}{%
      \text{\normalfont Schritte 1 und 5; (25.11) und Multiplikativität}}
    \proofstepwidestarbreak[1]{%
      s\mathrel{\sigma'_{\theta(\alpha),\theta(\beta)}}t
      \quad(s,t\in\Psi(x)\cup\Psi(y))
      \quad\text{\normalfont(25.20)}}{%
      \text{\normalfont Schritte 3--6; innere und bildübergreifende Äquivalenz}}
  \closeproofpartwide

  \proofpartwide{Transport der \(\kappa\)-Relation}
    \proofstepwidestar[1]{x\mathrel{\kappa_{\alpha,\gamma}}y}{\rA}
    \proofstepwidestar[2]{c'\in T_{\theta(\gamma)}}{\rA}
    \proofstepwidestar[2]{%
      A=\Psi^{-1}(\{c'\})\in\mathcal P_+(S_\gamma)}{%
      \text{\normalfont Komponententransport (25.4)}}
    \proofstepwidestar[2,4]{c\in A}{\rA}
    \proofstepwidestarbreak[1,2,4]{%
      AxA=\{cxc\}=\{cyc\}=AyA}{%
      \text{\normalfont Schritte 3--4; (25.10a) für \(\Psi^{-1}\) und (25.12)}}
    \proofstepwidestarbreak[1,2]{%
      c'\Psi(x)c'=c'\Psi(y)c'\text{ ist einelementig}}{%
      \text{\normalfont Schritt 5; Wahl von \(c\) eliminiert; dann (25.10b)}}
    \proofstepwidestarbreak[1]{%
      s\mathrel{\kappa'_{\theta(\alpha),\theta(\gamma)}}t
      \quad(s,t\in\Psi(x)\cup\Psi(y))
      \quad\text{\normalfont(25.21)}}{%
      \text{\normalfont Schritt 6; \(c'\) war beliebig; Definition (25.12)}}
  \closeproofpartwide

  \proofpartwide{Mengenfassungen des Transports}
    \proofstepwidestarbreak[1]{%
      \varnothing\neq D\subseteq a\sigma_{\alpha,\beta}\dsep
      u\in T_{\theta(\beta)}}{\rA}
    \proofstepwidestar[1]{%
      B=\Psi^{-1}(\{u\})\subseteq S_\beta\dsep B^2=B}{%
      \text{\normalfont (25.4), \(u^2=u\), Multiplikativität}}
    \proofstepwidestar[1]{DB=aB\dsep BD=Ba}{%
      \text{\normalfont Schritte 1--2; Definition (25.11)}}
    \proofstepwidestar[1]{DBD=aBa}{%
      \text{\normalfont Schritte 2--3; \(B^2=B\)}}
    \proofstepwidestarbreak[1]{%
      \Psi(D)u\Psi(D)=\Psi(a)u\Psi(a)
      \text{ ist einelementig}}{%
      \text{\normalfont Schritt 4; \(\Psi\) anwenden und Sandwichrechnung}}
    \proofstepwidestarbreak[1]{%
      s\mathrel{\sigma'_{\theta(\alpha),\theta(\beta)}}t
      \quad(s,t\in\Psi(a)\cup\Psi(D))
      \quad\text{\normalfont(25.20a)}}{%
      \text{\normalfont Schritt 5; Charakterisierung (25.13)}}
    \proofstepwidestarbreak[7]{%
      \varnothing\neq D\subseteq a\kappa_{\alpha,\gamma}\dsep
      v\in T_{\theta(\gamma)}}{\rA}
    \proofstepwidestar[7]{%
      A=\Psi^{-1}(\{v\})\subseteq S_\gamma}{%
      \text{\normalfont Komponententransport (25.4)}}
    \proofstepwidestarbreak[7]{%
      d\in D\cup\{a\}\vdash AdA\text{ ist einelementig}}{%
      \text{\normalfont Schritte 7--8; symmetrische Aussage für \(\Psi^{-1}\)}}
    \proofstepwidestar[7,10]{c\in A}{\rA}
    \proofstepwidestar[7,10]{AdA=\{cdc\}}{%
      \text{\normalfont Schritte 9--10; das Produkt enthält \(cdc\)}}
    \proofstepwidestar[7,10]{cdc=cac}{%
      \text{\normalfont Schritte 7 und 11; Definition (25.12)}}
    \proofstepwidestar[7]{ADA=AaA}{%
      \text{\normalfont Schritte 11--12; Wahl von \(c\) eliminiert}}
    \proofstepwidestarbreak[7]{%
      v\Psi(D)v=v\Psi(a)v\text{ ist einelementig}}{%
      \text{\normalfont Schritt 13; \(\Psi\) anwenden und (25.10b)}}
    \proofstepwidestarbreak[7]{%
      s\mathrel{\kappa'_{\theta(\alpha),\theta(\gamma)}}t
      \quad(s,t\in\Psi(a)\cup\Psi(D))
      \quad\text{\normalfont(25.21a)}}{%
      \text{\normalfont Schritt 14; Definition (25.12)}}
  \closeproofpartwide

  \proofpartwide{Erkennung und Transport der \(\rho\)-Klassen}
    \proofstepwidestar[1]{a\mathrel{\rho_\alpha}b}{\rA}
    \proofstepwidestarbreak[1]{%
      \begin{aligned}[t]
        &a\mathrel{\sigma_{\alpha,\beta}}b\quad(\beta<\alpha),\\[-2pt]
        &a\mathrel{\kappa_{\alpha,\gamma}}b\quad(\gamma>\alpha)
      \end{aligned}}{%
      \text{\normalfont Schnittdefinition (25.14)}}
    \proofstepwidestarbreak[1]{%
      \forall s,t\in\Psi(a)\cup\Psi(b)\;
      s\mathrel{\rho'_{\theta(\alpha)}}t}{%
      \text{\normalfont Schritt 2; (25.20), (25.21) und (25.14)}}
    \proofstepwidestarbreak[4]{%
      \forall s,t\in\Psi(a)\cup\Psi(b)\;
      s\mathrel{\rho'_{\theta(\alpha)}}t}{\rA}
    \proofstepwidestar[4,5]{s\in\Psi(a)\dsep
      A_s=\Psi^{-1}(\{s\})}{\rA}
    \proofstepwidestarbreak[4,5]{%
      \begin{aligned}[t]
        A_s\cup\Psi^{-1}(\Psi(a))&=A_s\cup\{a\},\\[-2pt]
        A_s\cup\Psi^{-1}(\Psi(b))&=A_s\cup\{b\}
      \end{aligned}
      \text{ liegen jeweils in einer }\rho_\alpha\text{-Klasse}}{%
      \text{\normalfont (25.20a), (25.21a) für \(\Psi^{-1}\)}}
    \proofstepwidestar[4]{a\mathrel{\rho_\alpha}b}{%
      \text{\normalfont Schritte 5--6; Wahl von \(s\) eliminiert; Brückenmenge \(A_s\)}}
    \proofstepwidestarbreak[]{%
      a\mathrel{\rho_\alpha}b
      \iff
      \forall s,t\in\Psi(a)\cup\Psi(b)\;
      s\mathrel{\rho'_{\theta(\alpha)}}t}{%
      \text{\normalfont Schritte 1--7; beide Annahmen entladen; dies ist (25.15)}}
    \proofstepwidestar[9]{C=a\rho_\alpha\dsep t\in\Psi(a)}{\rA}
    \proofstepwidestarbreak[9]{%
      \varnothing\neq D\subseteq C
      \vdash\Psi(D)\subseteq t\rho'_{\theta(\alpha)}}{%
      \text{\normalfont Schritte 8--9; Transport (25.20a), (25.21a)}}
    \proofstepwidestarbreak[9]{%
      \varnothing\neq E\subseteq t\rho'_{\theta(\alpha)}\dsep
      D=\Psi^{-1}(E)
      \vdash D\subseteq C}{%
      \text{\normalfont Schritte 8--9; inverser Transport mit \(\Psi^{-1}(\{t\})\)}}
    \proofstepwidestarbreak[9]{%
      \Psi\restriction_{\mathcal P_+(C)}\colon
      \mathcal P_+(C)\longrightarrow
      \mathcal P_+(t\rho'_{\theta(\alpha)})
      \text{ ist ein Isomorphismus}}{%
      \text{\normalfont Schritte 10--11; Bijektivität und Multiplikativität}}
  \closeproofpartwide

  \proofpartwide{Quotientenisomorphismus}
    \proofstepwidestar[]{%
      \tau_\alpha(a\rho_\alpha)
      =t\rho'_{\theta(\alpha)}\quad(t\in\Psi(a))}{%
      \text{\normalfont Vorschrift (25.16)}}
    \proofstepwidestar[]{\tau_\alpha\text{ ist wohldefiniert und injektiv}}{%
      \text{\normalfont Schritt 1 und Äquivalenz (25.15)}}
    \proofstepwidestar[3]{%
      E=t\rho'_{\theta(\alpha)}\text{ ist eine beliebige Zielklasse}}{\rA}
    \proofstepwidestar[3,4]{a\in\Psi^{-1}(\{t\})}{\rA}
    \proofstepwidestarbreak[3,4]{%
      \Psi^{-1}\restriction_{\mathcal P_+(E)}\colon
      \mathcal P_+(E)\bij\mathcal P_+(a\rho_\alpha)}{%
      \text{\normalfont Schritt 4; Klassenrestriktion für \(\Psi^{-1}\)}}
    \proofstepwidestar[3,4]{%
      \Psi(a)\subseteq E\dsep\tau_\alpha(a\rho_\alpha)=E}{%
      \text{\normalfont Schritte 1 und 5; inverse Klassenrestriktion}}
    \proofstepwidestar[]{\tau_\alpha\text{ ist bijektiv}}{%
      \text{\normalfont Schritte 2--6; Wahl von \(a\) und Zielklasse entladen}}
    \proofstepwidestarbreak[]{%
      a,b\in S_\alpha\dsep s\in\Psi(a)\dsep t\in\Psi(b)
      \vdash st\in\Psi(ab)}{%
      \text{\normalfont Multiplikativität von \(\Psi\)}}
    \proofstepwidestarbreak[]{%
      \tau_\alpha(a\rho_\alpha)\tau_\alpha(b\rho_\alpha)
      =(st)\rho'_{\theta(\alpha)}
      =\tau_\alpha((ab)\rho_\alpha)}{%
      \text{\normalfont Schritte 1 und 8; Quotientenprodukt}}
    \proofstepwidestar[]{%
      \tau_\alpha\colon S_\alpha/\rho_\alpha
      \longrightarrow T_{\theta(\alpha)}/\rho'_{\theta(\alpha)}
      \text{ ist ein Isomorphismus}}{%
      \text{\normalfont Schritte 7 und 9; Bijektivität und Multiplikativität; (25.17)}}
  \closeproofpartwide

  \proofpartwide{Produkt zweier echt tieferer Komponenten}
    \proofstepwidestar[1]{%
      \delta=\alpha\beta\dsep\alpha,\beta>\delta\dsep
      x\in a\rho_\alpha\dsep y\in b\rho_\beta}{\rA}
    \proofstepwidestar[1]{%
      x\mathrel{\sigma_{\alpha,\delta}}a\dsep
      y\mathrel{\sigma_{\beta,\delta}}b}{%
      \text{\normalfont Schnittdefinition (25.14)}}
    \proofstepwidestarbreak[1]{%
      \begin{aligned}[t]
        xy&=(xyx)y=(xyx)b=(xy)ab\\[-2pt]
          &=xy(ab)=xb(ab)=a(bab)=ab
      \end{aligned}
      \quad\text{\normalfont(25.22)}}{%
      \text{\normalfont Schritt 2; (25.11); Testelemente liegen in \(S_\delta\)}}
    \proofstepwidestarbreak[]{%
      (a\rho_\alpha)(b\rho_\beta)=\{ab\}}{%
      \text{\normalfont Schritt 3; \(x,y\) waren beliebige Vertreter}}
    \proofstepwidestarbreak[]{%
      \tau_\alpha(a\rho_\alpha)\text{ und }
      \tau_\beta(b\rho_\beta)
      \text{ liegen in den entsprechenden }\sigma'\text{-Klassen}}{%
      \text{\normalfont Definition (25.14)}}
    \proofstepwidestarbreak[]{%
      \tau_\alpha(a\rho_\alpha)\tau_\beta(b\rho_\beta)
      \text{ ist einelementig}}{%
      \text{\normalfont Schritte 3 und 5; Zielrechnung (25.22)}}
    \proofstepwidestarbreak[]{%
      s\in\Psi(a)\dsep t\in\Psi(b)
      \vdash st\in\Psi(ab)
      \cap\tau_\alpha(a\rho_\alpha)\tau_\beta(b\rho_\beta)}{%
      \text{\normalfont Schritt 6; Multiplikativität und Definition von \(\tau\)}}
    \proofstepwidestarbreak[]{%
      \Psi(ab)=\Psi(a)\Psi(b)\subseteq
      \tau_\alpha(a\rho_\alpha)\tau_\beta(b\rho_\beta)}{%
      \text{\normalfont Schritt 5; Multiplikativität und Definition (25.16)}}
    \proofstepwidestarbreak[]{%
      \tau_\alpha(a\rho_\alpha)\tau_\beta(b\rho_\beta)
      =\Psi(ab)\quad\text{\normalfont(25.18)}}{%
      \text{\normalfont Schritte 6--8; dasselbe einzige Element}}
  \closeproofpartwide

  \proofpartwide{Sandwichprodukt geordneter Komponenten}
    \proofstepwidestar[1]{%
      \alpha>\beta\dsep x\in a\rho_\alpha\dsep
      y\in b\rho_\beta}{\rA}
    \proofstepwidestar[1]{%
      x\mathrel{\sigma_{\alpha,\beta}}a\dsep
      y\mathrel{\kappa_{\beta,\alpha}}b}{%
      \text{\normalfont Schnittdefinition (25.14)}}
    \proofstepwidestar[1]{%
      xyx=aya=aba\quad\text{\normalfont(25.23)}}{%
      \text{\normalfont Schritt 2; Charakterisierungen (25.13), (25.12)}}
    \proofstepwidestarbreak[]{%
      \tau_\alpha(a\rho_\alpha)\tau_\beta(b\rho_\beta)
      \tau_\alpha(a\rho_\alpha)\text{ ist einelementig}}{%
      \text{\normalfont Schritt 3; Vertreter waren beliebig; Zielrechnung}}
    \proofstepwidestarbreak[]{%
      s\in\Psi(a)\dsep t\in\Psi(b)
      \vdash sts\in\Psi(aba)\cap
      \tau_\alpha(a\rho_\alpha)\tau_\beta(b\rho_\beta)
      \tau_\alpha(a\rho_\alpha)}{%
      \text{\normalfont Schritt 4; Multiplikativität und Definition von \(\tau\)}}
    \proofstepwidestar[]{\Psi(aba)\text{ ist einelementig}}{%
      \text{\normalfont Sandwicherkennung (25.5)}}
    \proofstepwidestarbreak[]{%
      \tau_\alpha(a\rho_\alpha)\tau_\beta(b\rho_\beta)
      \tau_\alpha(a\rho_\alpha)=\Psi(aba)
      \quad\text{\normalfont(25.19)}}{%
      \text{\normalfont Schritte 4--6; dasselbe einzige Element}}
  \closeproofpartwide
\end{tabproofsplitwide}
\endgroup

\section{Endliches Anheben auf die Grundelemente}

Der Quotientenisomorphismus \(\tau_\alpha\) muss nun innerhalb jeder
\(\rho_\alpha\)-Klasse zu einer Bijektion der Grundelemente angehoben
werden.  Bei echten Rechtecksbändern und bei trivialen Komponenten ist
diese Hebung bereits kanonisch.  Für echte Rechtecksbänder folgt die
Einermengeneigenschaft aus dem oben bewiesenen Erkennungssatz; bei einer
trivialen Komponente folgt sie unmittelbar aus (25.4), weil ihre
nichtleere Potenzmenge genau ein Element besitzt.  In beiden Fällen ist
\(\Psi(a)\) für jedes \(a\in S_\alpha\) eine Einermenge, und wir setzen
\[
  \eta_\alpha(a)=\text{das einzige Element von }\Psi(a).   \tag{25.24}
\]
Die so definierte Abbildung ist ein Isomorphismus
\(S_\alpha\to T_{\theta(\alpha)}\).  Wegen (25.16) liegt
\(\eta_\alpha(a)\) außerdem in der Klasse
\(\tau_\alpha(a\rho_\alpha)\).

Es bleiben Links- und Rechtsnullkomponenten.  Für eine solche Komponente
setzen wir
\[
U_\alpha=\begin{cases}
\displaystyle
 \bigcup_{\gamma>\alpha}\ \bigcup_{c\in S_\gamma}cS_\alpha,
   &\text{falls }S_\alpha\text{ linksnull ist},\\[3mm]
\displaystyle
 \bigcup_{\gamma>\alpha}\ \bigcup_{c\in S_\gamma}S_\alpha c,
   &\text{falls }S_\alpha\text{ rechtsnull ist}.
\end{cases}                                                \tag{25.25}
\]
Für eine zugeordnete Komponente von \(T\) bezeichnet
\(U_{\theta(\alpha)}\) im Folgenden die durch dieselbe Vorschrift in
\(T\) gebildete Menge.
Ist \(\alpha\) maximal, so ist \(U_\alpha=\varnothing\).  In beiden
Fällen gilt auch
\[
 U_\alpha=
 \bigcup_{\gamma>\alpha}\ \bigcup_{c\in S_\gamma}cS_\alpha c.  \tag{25.26}
\]
  Tatsächlich liegen \(cd\) und \(dc\) beide in \(S_\alpha\).  Ist diese
  Komponente linksnull, so ist
  \(cd=(cd)(dc)=cdc\); ist sie rechtsnull, so ist
  \(dc=(cd)(dc)=cdc\).  Dies gilt für
  \(c\in S_\gamma\), \(d\in S_\alpha\) und \(\gamma>\alpha\) und beweist
  (25.26).

\paragraph{Kanonischer Teil der Hebung.}
Für \(u\in U_\alpha\) ist \(\Psi(u)\) eine Einermenge, und ihr einziges
Element liegt in \(U_{\theta(\alpha)}\).  Die dadurch definierte Abbildung
\[
 U_\alpha\longrightarrow U_{\theta(\alpha)}              \tag{25.27}
\]
ist bijektiv.

\begingroup
\let\proofstepwidestar\proofstepwidestarforcebreak
\let\proofstepwidestarbreak\proofstepwidestarforcebreak
\ProofWideReasonsNearFormula
\begin{tabproofwide}
  \proofstepwidestar[1]{u\in U_\alpha}{\rA}
  \proofstepwidestarbreak[1]{%
    \exists\gamma>\alpha\,\exists c\in S_\gamma\,
      \exists d\in S_\alpha\;(u=cdc)}{%
    \text{\normalfont Darstellung (25.26)}}
  \proofstepwidestar[1]{\Psi(cdc)\text{ ist eine Einermenge}}{%
    \text{\normalfont Schritt 2; Sandwicherkennung (25.5)}}
  \proofstepwidestarbreak[1]{%
    \Psi(u)=\Psi(c)\Psi(d)\Psi(c)=\{sts\}
    \quad\text{für }s\in T_{\theta(\gamma)},\,
      t\in T_{\theta(\alpha)}}{%
    \text{\normalfont Schritte 2--3; Multiplikativität und (25.4)}}
  \proofstepwidestar[1]{sts\in U_{\theta(\alpha)}}{%
    \text{\normalfont Schritte 2 und 4; Darstellung (25.26) in \(T\)}}
  \proofstepwidestar[1]{%
    u\longmapsto\text{das einzige Element von }\Psi(u)
    \text{ definiert }U_\alpha\to U_{\theta(\alpha)}}{%
    \text{\normalfont Schritte 1 und 5; Wohldefiniertheit}}
  \proofstepwidestar[]{%
    v\longmapsto\text{das einzige Element von }\Psi^{-1}(\{v\})
    \text{ definiert }U_{\theta(\alpha)}\to U_\alpha}{%
    \text{\normalfont dieselbe Rechnung für \(\Psi^{-1}\)}}
  \proofstepwidestar[]{%
    \text{die beiden Zuordnungen sind zueinander invers}}{%
    \text{\normalfont Schritte 6--7; \(\Psi^{-1}\Psi=\Id\) und \(\Psi\Psi^{-1}=\Id\)}}
  \proofstepwidestar[]{U_\alpha\longrightarrow
    U_{\theta(\alpha)}\text{ ist bijektiv}}{%
    \text{\normalfont Schritte 6--8; inverse Abbildungen}}
\end{tabproofwide}
\endgroup

Sei jetzt \(C=a\rho_\alpha\) und
\[
 D=\tau_\alpha(C)
   =t\rho'_{\theta(\alpha)}\qquad(t\in\Psi(a)).            \tag{25.28}
\]
Nach dem Transportlemma schränkt sich \(\Psi\) zu einer Bijektion
\(\mathcal P_+(C)\to\mathcal P_+(D)\) ein.  Also sind die beiden
nichtleeren Potenzmengen gleichmächtig.  Weil \(C,D\) endlich sind,
liefert
\FormulaRefAuto{FiniteNonemptyPowerSetReflectsCardinality}[thm]
\[
  |C|=|D|.                                                  \tag{25.29}
\]
Aus (25.15) und (25.27) folgt außerdem, dass der kanonische Teil eine
Bijektion
\[
 C\cap U_\alpha\longrightarrow
 D\cap U_{\theta(\alpha)}                                 \tag{25.30}
\]
ist: Die Inklusion folgt aus (25.15), und die umgekehrte Inklusion folgt
durch dieselbe Argumentation für die inverse kanonische Abbildung.  Aus
(25.29) und (25.30) folgt, unter Verwendung der disjunkten Zerlegungen
\(C=(C\cap U_\alpha)\cup(C\setminus U_\alpha)\) und der
analogen Zerlegung von \(D\),
\[
\begin{aligned}
 |C\setminus U_\alpha|
 &=|C|-|C\cap U_\alpha|\\
 &=|D|-|D\cap U_{\theta(\alpha)}|\\
 &=|D\setminus U_{\theta(\alpha)}|.
\end{aligned}
\]
Ergänze (25.30) auf diesen Komplementen durch eine beliebige Bijektion.
So entsteht eine Bijektion
\[
  \overline\tau_C\colon C\longrightarrow D               \tag{25.31}
\]
mit
\[
 \overline\tau_C(u)=\text{das einzige Element von }\Psi(u)
 \qquad(u\in C\cap U_\alpha).                             \tag{25.32}
\]
Da \(\rho_\alpha\) und \(\rho'_{\theta(\alpha)}\) Kongruenzen sind und
die Grundelemente idempotent sind, sind \(C\) und \(D\) Unterhalbgruppen.
Jede Abbildung (25.31) ist ein Isomorphismus, denn Quelle und Ziel sind
beide Linksnull- beziehungsweise beide Rechtsnullhalbgruppen.

Die Menge der \(\rho_\alpha\)-Klassen ist das Bild der endlichen Menge
\(S_\alpha\) unter der Quotientenprojektion und daher endlich.  Durch
endlich viele aufeinanderfolgende Wahlen wählen wir für jede Klasse genau
eine solche Abbildung und vereinigen sie; wegen der paarweise disjunkten
Klassen ist diese Vereinigung eine Funktion.  Man erhält einen
Komponentenisomorphismus
\[
 \eta_\alpha\colon S_\alpha\longrightarrow
 T_{\theta(\alpha)}                                       \tag{25.33}
\]
mit den beiden Verträglichkeiten
\begin{align}
 \eta_\alpha(a)\rho'_{\theta(\alpha)}
    &=\tau_\alpha(a\rho_\alpha),                         \tag{25.34}\\
 a\in U_\alpha\quad&\Longrightarrow\quad
 \eta_\alpha(a)=\text{das einzige Element von }\Psi(a). \tag{25.35}
\end{align}
Dabei gilt (25.34) auch für die zuvor kanonisch definierten Abbildungen
aus (25.24).  Damit sind die Abbildungen \(\eta_\alpha\) für alle vier
Komponententypen definiert.

\section{Globale Bestimmtheit endlicher Bänder}

\FormulaThmDeltaK[Endliche idempotente Halbgruppen sind durch ihre
Potenzhalbgruppen global bestimmt]{%
  \begin{aligned}[t]
   &\IdempotentSemigroupAlg{A}{\star}\dsep
    \IdempotentSemigroupAlg{B}{\diamond}\dsep
    \Finite{A}\dsep\Finite{B}\dsep{}\\[-2pt]
   &\SemigroupIso{\Psi}
      {\PowNE{A},\star_{\mathcal P}}
      {\PowNE{B},\diamond_{\mathcal P}}
    \vdash
    \exists f\,\IdempotentSemigroupIso{f}{A,\star}{B,\diamond}
  \end{aligned}%
}{FiniteIdempotentSemigroupsGloballyDetermined}{%
  \DeltaRow{Mengen}{A\dsep B\dsep X\dsep Y}
  \DeltaRow{Funktionen}{f\dsep\Psi\dsep\star\dsep\diamond\dsep
    \star_{\mathcal P}\dsep\diamond_{\mathcal P}}
}

\begingroup
\let\proofstepwidestar\proofstepwidestarforcebreak
\let\proofstepwidestarbreak\proofstepwidestarforcebreak
\ProofWideReasonsNearFormula
\begin{tabproofsplitwide}
  \proofpartwide{Leerfall und Reduktion auf nichtleere Träger}
    \proofstepwidestarbreak[1]{%
      \begin{aligned}[t]
        &\IdempotentSemigroupAlg{A}{\star}\dsep
         \IdempotentSemigroupAlg{B}{\diamond}\dsep
         \Finite{A}\dsep\Finite{B},\\[-2pt]
        &\SemigroupIso{\Psi}
          {\PowNE{A},\star_{\mathcal P}}
          {\PowNE{B},\diamond_{\mathcal P}}
      \end{aligned}}{\rA}
    \proofstepwidestar[2]{A=\varnothing}{\rA}
    \proofstepwidestar[1,2]{\PowNE{A}=\varnothing}{%
      \text{Definition der nichtleeren Potenzmenge}}
    \proofstepwidestar[1,2]{\PowNE{B}=\varnothing}{%
      \text{Bijektivität von }\Psi\text{ und Schritt 3}}
    \proofstepwidestar[1,2]{B=\varnothing}{%
      \text{Definition der nichtleeren Potenzmenge}}
    \proofstepwidestarbreak[1,2]{%
      \exists f\,\IdempotentSemigroupIso{f}{A,\star}{B,\diamond}}{%
      \text{die leere Abbildung ist bijektiv und multiplikativ}}
    \proofstepwidestar[7]{A\neq\varnothing}{\rA}
    \proofstepwidestar[1,7]{\PowNE{A}\neq\varnothing}{%
      \text{eine Einermenge liegt in }\PowNE{A}}
    \proofstepwidestar[1,7]{\PowNE{B}\neq\varnothing}{%
      \text{Bijektivität von }\Psi\text{ und Schritt 8}}
    \proofstepwidestar[1,7]{B\neq\varnothing}{%
      \text{Definition der nichtleeren Potenzmenge}}
  \closeproofpartwide

  \proofpartwide{Komponentenweise Konstruktion des Kandidaten}
    \proofstepwidestarbreak[1]{%
      \begin{aligned}[t]
        &\IdempotentSemigroupAlg{A}{\star}\dsep
         \IdempotentSemigroupAlg{B}{\diamond}\dsep
         \Finite{A}\dsep\Finite{B},\\[-2pt]
        &\SemigroupIso{\Psi}
          {\PowNE{A},\star_{\mathcal P}}
          {\PowNE{B},\diamond_{\mathcal P}},\\[-2pt]
        &A,B\neq\varnothing
      \end{aligned}}{\rA}
    \proofstepwidestar[1]{S=(A,\star)\dsep T=(B,\diamond)}{%
      \text{Bezeichnung}}
    \proofstepwidestarbreak[1]{%
      S=[Y;S_\alpha]\dsep T=[Z;T_\gamma]}{%
      \text{McLean-Zerlegung (25.2)}}
    \proofstepwidestarbreak[1]{%
      \theta\colon Y\bij Z\dsep
      \Psi[\mathcal P_+(S_\alpha)]
        =\mathcal P_+(T_{\theta(\alpha)})}{%
      \text{Komponentensatz (25.4)}}
    \proofstepwidestarbreak[1]{%
      \forall\alpha\in Y\ \exists\eta_\alpha\colon
      S_\alpha\bij T_{\theta(\alpha)}\quad
      \text{mit (25.34)--(25.35)}}{%
      \text{endliches Anheben (25.24)--(25.35)}}
    \proofstepwidestar[1]{%
      \eta=\displaystyle\bigcup_{\alpha\in Y}\eta_\alpha
      \quad\text{\normalfont(25.36)}}{%
      \text{Definition}}
    \proofstepwidestarbreak[1]{\eta\colon S\bij T}{%
      \text{disjunkte Komponenten, Bijektivität von }\theta
      \text{ und der }\eta_\alpha}
  \closeproofpartwide

  \proofpartwide{Fall 1: Beide Faktoren liegen in derselben Komponente}
    \proofstepwidestarbreak[1]{%
      a\in S_\alpha\dsep b\in S_\beta\dsep\alpha=\beta}{\rA}
    \proofstepwidestar[1]{ab\in S_\alpha}{\text{(25.3)}}
    \proofstepwide[1]{\eta(ab)}{=}{\eta_\alpha(ab)}{\text{(25.36)}}
    \proofstepwide[1]{\eta_\alpha(ab)}{=}{%
      \eta_\alpha(a)\eta_\alpha(b)}{\text{(25.33)}}
    \proofstepwide[1]{\eta_\alpha(a)\eta_\alpha(b)}{=}{%
      \eta(a)\eta(b)}{\text{(25.36)}}
    \proofstepwidestar[1]{%
      \eta(ab)=\eta(a)\eta(b)\quad\text{\normalfont(25.37)}}{%
      \rChain{3,4,5}}
  \closeproofpartwide

  \proofpartwide{Fall 2: Das Produkt liegt echt unter beiden Komponenten}
    \proofstepwidestarbreak[1]{%
      a\in S_\alpha\dsep b\in S_\beta\dsep
      \delta=\alpha\beta\notin\{\alpha,\beta\}}{\rA}
    \proofstepwidestar[1]{\alpha,\beta>\delta\dsep ab\in S_\delta}{%
      \text{Halbverbandsordnung und (25.3)}}
    \proofstepwidestarbreak[1]{%
      S_\delta\text{ trivial oder echtes Rechtecksband}
      \Longrightarrow
      \eta(ab)\text{ ist das einzige Element von }\Psi(ab)}{%
      \text{kanonische Hebung (25.24)}}
    \proofstepwidestarbreak[1]{%
      S_\delta\text{ linksnull}
      \Longrightarrow ab=a(ab)\in U_\delta}{%
      \text{(25.25) und Linksnullidentität}}
    \proofstepwidestarbreak[1]{%
      S_\delta\text{ rechtsnull}
      \Longrightarrow ab=(ab)b\in U_\delta}{%
      \text{(25.25) und Rechtsnullidentität}}
    \proofstepwidestarbreak[1]{%
      \eta(ab)\text{ ist das einzige Element von }\Psi(ab)}{%
      \text{Schritte 3--5 und (25.35)}}
    \proofstepwidestarbreak[1]{%
      \eta(a)\in\tau_\alpha(a\rho_\alpha)\dsep
      \eta(b)\in\tau_\beta(b\rho_\beta)}{\text{(25.34)}}
    \proofstepwidestarbreak[1]{%
      \eta(a)\eta(b)\text{ ist das einzige Element von }\Psi(ab)}{%
      \text{Produkttransport (25.18)}}
    \proofstepwidestar[1]{%
      \eta(ab)=\eta(a)\eta(b)\quad\text{\normalfont(25.38)}}{%
      \text{Schritte 6 und 8}}
  \closeproofpartwide

  \proofpartwide{Fall 3a: Echte Rechtecks- oder triviale Unterkomponente}
    \proofstepwidestarbreak[1]{%
      a\in S_\alpha\dsep b\in S_\beta\dsep
      \alpha=\alpha\beta<\beta\dsep
      S_\alpha\text{ trivial oder echtes Rechtecksband}}{\rA}
    \proofstepwidestarbreak[1]{%
      \Psi(a)=\{\eta(a)\}\dsep
      \Psi(ab)=\{\eta(ab)\}}{\text{kanonische Hebung (25.24)}}
    \proofstepwidestarbreak[1]{%
      \exists C\in\mathcal P_+(b\rho_\beta)\quad
      \Psi(C)=\{\eta(b)\}}{\text{Klassenrestriktion des Transportlemmas}}
    \proofstepwidestar[1]{C\subseteq b\sigma_{\beta,\alpha}}{%
      \rho_\beta\subseteq\sigma_{\beta,\alpha}\text{ nach (25.14)}}
    \proofstepwidestarbreak[1]{%
      \forall s\in\Psi(b)\quad
      \eta(b)\mathrel{\sigma'_{\theta(\beta),\theta(\alpha)}}s}{%
      \text{mengenweiser Transport (25.20a)}}
    \proofstepwidestarbreak[1]{%
      \Psi(a)\Psi(C)=\Psi(a)\{s\}=\Psi(a)\Psi(b)
      \quad\text{\normalfont(25.39)}}{%
      \text{(25.11) im Ziel}}
    \proofstepwide[1]{\eta(a)\eta(b)}{=}{\Psi(a)\Psi(C)}{%
      \text{Schritte 2--3}}
    \proofstepwide[1]{\Psi(a)\Psi(C)}{=}{\Psi(ab)}{%
      \text{Multiplikativität von }\Psi}
    \proofstepwide[1]{\Psi(ab)}{=}{\eta(ab)}{\text{Schritt 2}}
    \proofstepwidestar[1]{%
      \eta(a)\eta(b)=\eta(ab)\quad\text{\normalfont(25.40)}}{%
      \rChain{7,8,9}}
  \closeproofpartwide

  \proofpartwide{Fall 3b: Linksnull-Unterkomponente}
    \proofstepwidestarbreak[1]{%
      a\in S_\alpha\dsep b\in S_\beta\dsep
      \alpha=\alpha\beta<\beta\dsep S_\alpha\text{ linksnull}}{\rA}
    \proofstepwide[1]{ab}{=}{a(ab)}{%
      \text{Assoziativität und }\(a^2=a\)}
    \proofstepwide[1]{a(ab)}{=}{a}{\text{Linksnullidentität}}
    \proofstepwidestar[1]{ab=a}{\rChain{2,3}}
    \proofstepwidestarbreak[1]{%
      x=\eta(a)\eta(b)\in T_{\theta(\alpha)}}{\text{(25.3)--(25.4)}}
    \proofstepwide[1]{\eta(a)x}{=}{\eta(a)}{%
      \text{Linksnullidentität im Ziel}}
    \proofstepwide[1]{\eta(a)x}{=}{%
      \eta(a)^2\eta(b)=x}{\text{Definition von }x\text{ und Idempotenz}}
    \proofstepwidestar[1]{%
      \eta(a)\eta(b)=\eta(a)=\eta(ab)
      \quad\text{\normalfont(25.41)}}{%
      \text{Schritte 4--7}}
  \closeproofpartwide

  \proofpartwide{Fall 3c: Rechtsnull-Unterkomponente}
    \proofstepwidestarbreak[1]{%
      a\in S_\alpha\dsep b\in S_\beta\dsep
      \alpha=\alpha\beta<\beta\dsep S_\alpha\text{ rechtsnull}}{\rA}
    \proofstepwidestar[1]{ba,ab\in S_\alpha}{\text{(25.3)}}
    \proofstepwide[1]{ab}{=}{(ba)(ab)}{\text{Rechtsnullidentität}}
    \proofstepwide[1]{(ba)(ab)}{=}{bab}{\text{Assoziativität}}
    \proofstepwidestar[1]{%
      ab=bab\in U_\alpha\quad\text{\normalfont(25.42)}}{%
      \text{Schritte 2--4 und (25.25)}}
    \proofstepwidestarbreak[1]{%
      \eta(ab)\text{ ist das einzige Element von }\Psi(bab)}{%
      \text{Schritt 5 und (25.35)}}
    \proofstepwidestarbreak[1]{%
      x=\eta(a)\eta(b),\ y=\eta(b)\eta(a)
      \in T_{\theta(\alpha)}}{\text{(25.3)--(25.4)}}
    \proofstepwide[1]{\eta(a)\eta(b)}{=}{%
      \eta(b)\eta(a)\eta(b)\quad\text{\normalfont(25.43)}}{%
      \text{Rechtsnullidentität }\(yx=x\)}
    \proofstepwidestarbreak[1]{%
      \eta(b)\eta(a)\eta(b)\in
      \tau_\beta(b\rho_\beta)\tau_\alpha(a\rho_\alpha)
      \tau_\beta(b\rho_\beta)}{\text{(25.34)}}
    \proofstepwidestarbreak[1]{%
      \tau_\beta(b\rho_\beta)\tau_\alpha(a\rho_\alpha)
      \tau_\beta(b\rho_\beta)=\Psi(bab)}{%
      \text{Sandwichtransport (25.19), da }\beta>\alpha}
    \proofstepwidestar[1]{\eta(a)\eta(b)=\eta(ab)}{%
      \text{Schritte 6 und 8--10}}
  \closeproofpartwide

  \proofpartwide{Fall 3d: Die umgekehrte Vergleichsrichtung}
    \proofstepwidestarbreak[1]{%
      a\in S_\alpha\dsep b\in S_\beta\dsep
      \beta=\alpha\beta<\alpha}{\rA}
    \proofstepwidestarbreak[1]{%
      \Psi\colon\mathcal P_+(S^{\mathrm{op}})
        \longrightarrow\mathcal P_+(T^{\mathrm{op}})
      \text{ ist ein Isomorphismus}}{%
      \text{beide Potenzhalbgruppenprodukte werden umgekehrt}}
    \proofstepwidestarbreak[1]{%
      \sigma,\kappa,\rho,\tau\text{ bleiben erhalten, }
      U_\gamma\text{ bleibt dieselbe Menge}}{%
      \text{Definitionen (25.11)--(25.16) und (25.25)}}
    \proofstepwidestarbreak[1]{%
      b\mathbin{\circ}a=ab\dsep
      \eta(b)\mathbin{\circ}\eta(a)=\eta(a)\eta(b)}{%
      \text{Definition der Oppositionsmultiplikation}}
    \proofstepwidestar[1]{\eta(a)\eta(b)=\eta(ab)}{%
      \text{Fälle 3a--3c für }(b,a)\text{ im Oppositum}}
  \closeproofpartwide

  \proofpartwide{Fallunterscheidung und Schluss}
    \proofstepwidestarbreak[1]{%
      \begin{aligned}[t]
        &\IdempotentSemigroupAlg{A}{\star}\dsep
         \IdempotentSemigroupAlg{B}{\diamond}\dsep
         \Finite{A}\dsep\Finite{B},\\[-2pt]
        &\SemigroupIso{\Psi}
          {\PowNE{A},\star_{\mathcal P}}
          {\PowNE{B},\diamond_{\mathcal P}}
      \end{aligned}}{\rA}
    \proofstepwidestar[2]{A\neq\varnothing}{\rA}
    \proofstepwidestar[1,2]{B\neq\varnothing}{%
      \text{Reduktion im ersten Teil}}
    \proofstepwidestarbreak[1,2]{%
      \eta=\displaystyle\bigcup_{\alpha\in Y}\eta_\alpha
      \colon S\bij T}{%
      \text{komponentenweise Konstruktion}}
    \proofstepwidestarbreak[5]{%
      a\in S_\alpha\dsep b\in S_\beta\dsep
      \delta=\alpha\beta\leq\alpha,\beta}{\rA}
    \proofstepwidestarbreak[1,2,5]{%
      \alpha=\beta\ \lor\
      \delta\notin\{\alpha,\beta\}\ \lor\
      \delta=\alpha<\beta\ \lor\
      \delta=\beta<\alpha}{\text{erschöpfende Lage der beiden Indizes}}
    \proofstepwidestar[1,2,5]{\eta(ab)=\eta(a)\eta(b)}{%
      \text{Fälle 1, 2 und 3a--3d}}
    \proofstepwidestarbreak[1,2]{%
      \forall a,b\in S\quad\eta(ab)=\eta(a)\eta(b)}{%
      \text{Allgemeinheit von }a,b\text{; Schritte 5--7}}
    \proofstepwidestar[1,2]{\SemigroupIso{\eta}{S}{T}}{%
      \text{Bijektivität aus Schritt 4 und Multiplikativität aus Schritt 8}}
    \proofstepwidestarbreak[1,2]{%
      \IdempotentSemigroupIso{\eta}{A,\star}{B,\diamond}}{%
      \FormulaRefAuto{IdempotentSemigroupIsoDef}[def]}
    \proofstepwidestarbreak[1,2]{%
      \exists f\,\IdempotentSemigroupIso{f}{A,\star}{B,\diamond}}{%
      \text{Existenz-Einführung mit }f=\eta}
    \proofstepwidestarbreak[1]{%
      A=\varnothing\Longrightarrow
      \exists f\,\IdempotentSemigroupIso{f}{A,\star}{B,\diamond}}{%
      \text{Leerfall im ersten Teil}}
    \proofstepwidestarbreak[1]{%
      \exists f\,\IdempotentSemigroupIso{f}{A,\star}{B,\diamond}}{%
      \text{Fallunterscheidung }A=\varnothing\text{ bzw. }A\neq\varnothing;
      \text{ Schritte 11--12}}
  \closeproofpartwide
\end{tabproofsplitwide}
\endgroup

\section{Reichweite der Aussage}

\begin{remark}
Der Satz setzt auf beiden Seiten Idempotenz voraus.  Der Beweis erkennt
nicht aus \(\mathcal P_+(S)\cong\mathcal P_+(T)\), dass eine zunächst
beliebige Zielhalbgruppe \(T\) ebenfalls idempotent ist.  Eine solche
einseitige Erkennung wäre ein zusätzlicher und für die allgemeine
Potenzhalbgruppenvermutung wesentlicher Schritt.

Die in
\href{https://doi.org/10.1080/00927872.2017.1319474}%
 {Yu--Zhao--Gan, \emph{Global determinism of idempotent semigroups}, 2018}
ohne Endlichkeitsbedingung formulierte Aussage wird hier bewusst nicht
übernommen.  Für ein Linksnullband \(L_X\) ist
\(\mathcal P_+(L_X)\) selbst ein Linksnullband mit
\(2^{|X|}-1\) Elementen.  Bei unendlichen Kardinalzahlen folgt aus
\(2^\kappa=2^\lambda\) in ZFC nicht \(\kappa=\lambda\); entsprechende
Kontinuumsfunktionen sind nach
\href{https://doi.org/10.1016/0003-4843(70)90012-4}%
 {Eastons Satz}
mit ZFC verträglich.  Im endlichen Fall schließt dagegen
\FormulaRefAuto{FiniteNonemptyPowerSetReflectsCardinality}[thm]
genau diese Lücke.
\end{remark}

\end{document}
